\documentclass[11pt]{article}
\usepackage{fontspec,amsmath,amssymb,amsthm,hyperref}
\usepackage[margin=1in]{geometry}
\setlength{\emergencystretch}{3em}
\title{Concrete independent initialization and Brownian paths}
\author{LangevinLean mathematical exposition}
\date{}
\begin{document}\maketitle
\noindent Source module: \texttt{ProximalLangevin/MeanFieldInitialSource.lean}. Compile this proof document with XeLaTeX or LuaLaTeX. The accompanying HTML places the complete Lean declarations beside every proof.


Given a latent probability law \(\mu\) with finite second moment and a deterministic \(\theta_0\), use the product \(Q=\mu\otimes Q_B\), where \(Q_B\) is the canonical joint Brownian source. This constructs all initialization and source hypotheses from primitive data.


\section{Product source and its measurable Brownian evaluations}


The source is the actual product law \(\mu\otimes Q_B\). Brownian coordinates use only its second factor, so measure-preserving second projection transfers their standard Brownian laws. Coordinate measurability assembles into full-path measurability. The joint path pushforward remains exactly its canonical product law, proving independence of the complete coordinate paths after initialization is added.



The named probability instance follows from the product measure identity \(Q(\Omega)=\mu(X)Q_B(\Omega_B)=1\cdot1=1\). This records normalization of the constructed source before any law or expectation is used.



\noindent Lean declarations: \texttt{meanFieldInitialSource}, \texttt{meanFieldInitialBrownian}, \texttt{meanFieldInitialBrownian\_measurable}, \texttt{meanFieldInitialBrownian\_preBrownian}, \texttt{meanFieldInitialBrownian\_paths\_independent}, \texttt{meanFieldInitialSource\_isProbability}.

\section{Build the square-integrable pair initialization}


The map \(x\mapsto(\theta_0,x)\) is nonexpansive because its parameter difference vanishes. The latent first-factor variable has finite second moment under the product law. Lipschitz composition therefore proves that \((\theta_0,X_0)\) is square integrable, and its actual \(L^2\) representative defines the initialization.



\noindent Lean declarations: \texttt{meanFieldInitialPair\_lipschitz}, \texttt{meanFieldInitial\_memLp}, \texttt{meanFieldInitialLp}.

\section{Retain the exact parameter, latent law and source independence}


The representative equals \((\theta_0,\omega_1)\) almost everywhere, so its parameter projection is \(\operatorname{const}\theta_0\), and its latent pushforward is exactly \(\mu\). The whole initialization depends only on the first factor while all Brownian paths depend only on the second. Product independence, followed by the almost everywhere representative substitution, proves independence of initialization and the entire Brownian family.



\noindent Lean declarations: \texttt{meanFieldInitialLp\_parameter}, \texttt{meanFieldInitialLp\_latent\_law}, \texttt{meanFieldInitialLp\_independent\_paths}.

\end{document}
